% --- Archon physics formalization source begin ---
% archon:physics
% archon:covers PhyXMiniProblems/problem_phyx_mini_0149.lean
% archon:source-report reports/phyx_mini/problem_phyx_mini_0149.source.json

\chapter{Physics problem phyx\_mini\_0149}
\label{ch:PhyXMiniProblems_problem_phyx_mini_0149}

\paragraph{Problem source.}
\#\# Physical scenario

Figure shows a converging lens held above three equal-sized letters A. In (a) the lens is \textbackslash{}(5\textbackslash{},\textbackslash{}mathrm\{cm\}\textbackslash{}) from the paper, and in (b) the lens is \textbackslash{}(15\textbackslash{},\textbackslash{}mathrm\{cm\}\textbackslash{}) from the paper.

\#\# Question

Estimate the focal length of the lens.

\#\# Answer choices

- A: A: 8cm

- B: B: 12cm

- C: C: 10cm

- D: D: 14cm

\#\# Figure caption (auxiliary; use the image as primary evidence)

The image contains two parts, labeled (a) and (b), showing a metallic concave mirror reflecting the letter "A."

1. **Image (a):**
   - A concave mirror is positioned vertically, reflecting the letter "A."
   - The letter "A" is written on a lined paper in the background.
   - On either side of the mirror, the letter "A" appears on the paper.
   - The reflection of the "A" inside the mirror appears upright and slightly magnified.

2. **Image (b):**
   - The same concave mirror now shows an inverted reflection of the letter "A."
   - The letter "A" is again written on the same lined paper background and flanked by additional "A" letters.
   - The reflection of "A" in the mirror is upside down, indicating a change in the positioning relative to the focal point of the mirror.

The images demonstrate the reflecting properties of a concave mirror at different distances from the object.

\#\# Dataset metadata

Geometrical Optics; Physical Model Grounding Reasoning, Spatial Relation Reasoning

\paragraph{Recorded answer/context.}
C: C: 10cm

\paragraph{Figure/image path.}
/root/proposal\_for\_physic/science-mango/phyx\_mini\_run/phyx\_data/test\_image/149.png

\paragraph{Formalization target.}
create a compiling Lean file with sorry bodies at `PhyXMiniProblems/problem\_phyx\_mini\_0149.lean`. The Lean declarations must preserve the physical quantities, dimensions or dimensional roles, figure labels, governing-law hypotheses, and final relation expressed by this problem.
Use Mathlib/Physlib names found through LeanExplore where available. If a physics API is missing, introduce faithful local abstractions rather than scalar placeholder aliases.

\begin{theorem}[Physics formalization target]
\label{thm:physics:phyx_mini_0149:target}
\lean{PhyXMiniProblems.ProblemPhyXMini0149.problem_phyx_mini_0149}
\uses{def:physics:phyx-mini-0149:phyxminiproblems-problemphyxmini0149-lengthincentimeters, def:physics:phyx-mini-0149:phyxminiproblems-problemphyxmini0149-focallengthestimatesetup, def:physics:phyx-mini-0149:phyxminiproblems-problemphyxmini0149-matchesproblemandprimaryfigure, def:physics:phyx-mini-0149:phyxminiproblems-problemphyxmini0149-haspositivephysicallengths, def:physics:phyx-mini-0149:phyxminiproblems-problemphyxmini0149-satisfiesparaxialconverginglensmodel, def:physics:phyx-mini-0149:phyxminiproblems-problemphyxmini0149-isuniquematchinganswerchoice}
The assigned autoformalize agent should translate this physics problem into Lean declarations in the covered file, with theorem and lemma proof bodies written as `by sorry`.
\end{theorem}
\begin{proof}
This is an autoformalization task, not a proof task. Produce statements that can later be proved by the physics prover without weakening the physical model.
\end{proof}

% --- Archon declaration topology begin ---
\section{Declaration topology}
\begin{definition}[Length Quantity]
  \label{def:physics:phyx-mini-0149:phyxminiproblems-problemphyxmini0149-lengthquantity}
  \lean{PhyXMiniProblems.ProblemPhyXMini0149.LengthQuantity}
  A nonnegative physical length, independent of the unit used to read it
\end{definition}
\begin{proof}
  This is the declared model definition of Length Quantity.
\end{proof}
\begin{definition}[length In Centimeters]
  \label{def:physics:phyx-mini-0149:phyxminiproblems-problemphyxmini0149-lengthincentimeters}
  \lean{PhyXMiniProblems.ProblemPhyXMini0149.lengthInCentimeters}
  \uses{def:physics:phyx-mini-0149:phyxminiproblems-problemphyxmini0149-lengthquantity}
  The scalar readout of a physical length in centimeters
\end{definition}
\begin{proof}
  This is the declared model definition of length In Centimeters.
\end{proof}
\begin{definition}[Figure Panel]
  \label{def:physics:phyx-mini-0149:phyxminiproblems-problemphyxmini0149-figurepanel}
  \lean{PhyXMiniProblems.ProblemPhyXMini0149.FigurePanel}
  Labels (a) and (b) of the two source-image panels
\end{definition}
\begin{proof}
  This is the declared model definition of Figure Panel.
\end{proof}
\begin{definition}[Letter Position]
  \label{def:physics:phyx-mini-0149:phyxminiproblems-problemphyxmini0149-letterposition}
  \lean{PhyXMiniProblems.ProblemPhyXMini0149.LetterPosition}
  Horizontal positions of the three equal-sized printed letters
\end{definition}
\begin{proof}
  This is the declared model definition of Letter Position.
\end{proof}
\begin{definition}[Figure Element]
  \label{def:physics:phyx-mini-0149:phyxminiproblems-problemphyxmini0149-figureelement}
  \lean{PhyXMiniProblems.ProblemPhyXMini0149.FigureElement}
  \uses{def:physics:phyx-mini-0149:phyxminiproblems-problemphyxmini0149-letterposition}
  Physical elements present in each panel of the photograph
\end{definition}
\begin{proof}
  This is the declared model definition of Figure Element.
\end{proof}
\begin{definition}[Printed Glyph]
  \label{def:physics:phyx-mini-0149:phyxminiproblems-problemphyxmini0149-printedglyph}
  \lean{PhyXMiniProblems.ProblemPhyXMini0149.PrintedGlyph}
  The printed glyph used at all three positions
\end{definition}
\begin{proof}
  This is the declared model definition of Printed Glyph.
\end{proof}
\begin{definition}[Lens Kind]
  \label{def:physics:phyx-mini-0149:phyxminiproblems-problemphyxmini0149-lenskind}
  \lean{PhyXMiniProblems.ProblemPhyXMini0149.LensKind}
  Optical kind of a paraxial thin lens
\end{definition}
\begin{proof}
  This is the declared model definition of Lens Kind.
\end{proof}
\begin{definition}[Image Orientation]
  \label{def:physics:phyx-mini-0149:phyxminiproblems-problemphyxmini0149-imageorientation}
  \lean{PhyXMiniProblems.ProblemPhyXMini0149.ImageOrientation}
  Orientation of the central letter as viewed through the lens
\end{definition}
\begin{proof}
  This is the declared model definition of Image Orientation.
\end{proof}
\begin{definition}[Apparent Size Relation]
  \label{def:physics:phyx-mini-0149:phyxminiproblems-problemphyxmini0149-apparentsizerelation}
  \lean{PhyXMiniProblems.ProblemPhyXMini0149.ApparentSizeRelation}
  Apparent size of the central image relative to the uncovered letters
\end{definition}
\begin{proof}
  This is the declared model definition of Apparent Size Relation.
\end{proof}
\begin{definition}[Thin Lens]
  \label{def:physics:phyx-mini-0149:phyxminiproblems-problemphyxmini0149-thinlens}
  \lean{PhyXMiniProblems.ProblemPhyXMini0149.ThinLens}
  \uses{def:physics:phyx-mini-0149:phyxminiproblems-problemphyxmini0149-lengthquantity, def:physics:phyx-mini-0149:phyxminiproblems-problemphyxmini0149-lenskind}
  A thin lens with its genuine, as-yet unknown physical focal length
\end{definition}
\begin{proof}
  This is the declared model definition of Thin Lens.
\end{proof}
\begin{definition}[Focal Length Estimate Setup]
  \label{def:physics:phyx-mini-0149:phyxminiproblems-problemphyxmini0149-focallengthestimatesetup}
  \lean{PhyXMiniProblems.ProblemPhyXMini0149.FocalLengthEstimateSetup}
  \uses{def:physics:phyx-mini-0149:phyxminiproblems-problemphyxmini0149-lengthquantity, def:physics:phyx-mini-0149:phyxminiproblems-problemphyxmini0149-figurepanel, def:physics:phyx-mini-0149:phyxminiproblems-problemphyxmini0149-letterposition, def:physics:phyx-mini-0149:phyxminiproblems-problemphyxmini0149-figureelement, def:physics:phyx-mini-0149:phyxminiproblems-problemphyxmini0149-printedglyph, def:physics:phyx-mini-0149:phyxminiproblems-problemphyxmini0149-imageorientation, def:physics:phyx-mini-0149:phyxminiproblems-problemphyxmini0149-apparentsizerelation, def:physics:phyx-mini-0149:phyxminiproblems-problemphyxmini0149-thinlens}
  The physical setup and all panel-indexed measurements. The actual focal length remains an unknown physical length; no numerical answer is assigned to it here. Signed transverse magnification is positive for an upright image and negative for an inverted image
\end{definition}
\begin{proof}
  This is the declared model definition of Focal Length Estimate Setup.
\end{proof}
\begin{definition}[Matches Problem And Primary Figure]
  \label{def:physics:phyx-mini-0149:phyxminiproblems-problemphyxmini0149-matchesproblemandprimaryfigure}
  \lean{PhyXMiniProblems.ProblemPhyXMini0149.MatchesProblemAndPrimaryFigure}
  \uses{def:physics:phyx-mini-0149:phyxminiproblems-problemphyxmini0149-lengthincentimeters, def:physics:phyx-mini-0149:phyxminiproblems-problemphyxmini0149-focallengthestimatesetup}
  Primary-figure and problem-statement readouts: both panels contain the lens, paper, and three copies of A; the printed copies have equal physical height; the lens-paper separations are 5 cm and 15 cm; and the viewed central A is upright in (a) and inverted in (b). Both central images are magnified by approximately the same factor relative to the uncovered equal-sized letters; that visual estimate is idealized as equality of magnification magnitudes. No focal-length value or answer choice occurs in this predicate
\end{definition}
\begin{proof}
  This is the declared model definition of Matches Problem And Primary Figure.
\end{proof}
\begin{definition}[Has Positive Physical Lengths]
  \label{def:physics:phyx-mini-0149:phyxminiproblems-problemphyxmini0149-haspositivephysicallengths}
  \lean{PhyXMiniProblems.ProblemPhyXMini0149.HasPositivePhysicalLengths}
  \uses{def:physics:phyx-mini-0149:phyxminiproblems-problemphyxmini0149-lengthincentimeters, def:physics:phyx-mini-0149:phyxminiproblems-problemphyxmini0149-focallengthestimatesetup}
  Positivity conditions for every physical length used in the model
\end{definition}
\begin{proof}
  This is the declared model definition of Has Positive Physical Lengths.
\end{proof}
\begin{definition}[Satisfies Paraxial Converging Lens Model]
  \label{def:physics:phyx-mini-0149:phyxminiproblems-problemphyxmini0149-satisfiesparaxialconverginglensmodel}
  \lean{PhyXMiniProblems.ProblemPhyXMini0149.SatisfiesParaxialConvergingLensModel}
  \uses{def:physics:phyx-mini-0149:phyxminiproblems-problemphyxmini0149-lengthincentimeters, def:physics:phyx-mini-0149:phyxminiproblems-problemphyxmini0149-focallengthestimatesetup}
  The paraxial laws used in this problem. For a converging lens and a real object on the paper, an upright view means that the object is inside the focal plane and has positive signed magnification, whereas an inverted view means that it is beyond the focal plane and has negative signed magnification. The transverse-magnification equation is written without division as m * (f - u) = f, where all length readouts use centimeters.
\end{definition}
\begin{proof}
  This is the declared model definition of Satisfies Paraxial Converging Lens Model.
\end{proof}
\begin{definition}[Answer Choice]
  \label{def:physics:phyx-mini-0149:phyxminiproblems-problemphyxmini0149-answerchoice}
  \lean{PhyXMiniProblems.ProblemPhyXMini0149.AnswerChoice}
  Labels of the four focal-length choices printed in the problem
\end{definition}
\begin{proof}
  This is the declared model definition of Answer Choice.
\end{proof}
\begin{definition}[displayed Focal Length In Centimeters]
  \label{def:physics:phyx-mini-0149:phyxminiproblems-problemphyxmini0149-displayedfocallengthincentimeters}
  \lean{PhyXMiniProblems.ProblemPhyXMini0149.displayedFocalLengthInCentimeters}
  \uses{def:physics:phyx-mini-0149:phyxminiproblems-problemphyxmini0149-answerchoice}
  The focal-length readout printed beside each answer label, in centimeters
\end{definition}
\begin{proof}
  This is the declared model definition of displayed Focal Length In Centimeters.
\end{proof}
\begin{definition}[Matches Answer Choice]
  \label{def:physics:phyx-mini-0149:phyxminiproblems-problemphyxmini0149-matchesanswerchoice}
  \lean{PhyXMiniProblems.ProblemPhyXMini0149.MatchesAnswerChoice}
  \uses{def:physics:phyx-mini-0149:phyxminiproblems-problemphyxmini0149-lengthincentimeters, def:physics:phyx-mini-0149:phyxminiproblems-problemphyxmini0149-focallengthestimatesetup, def:physics:phyx-mini-0149:phyxminiproblems-problemphyxmini0149-answerchoice, def:physics:phyx-mini-0149:phyxminiproblems-problemphyxmini0149-displayedfocallengthincentimeters}
  Exact agreement between the inferred physical focal length and a choice
\end{definition}
\begin{proof}
  This is the declared model definition of Matches Answer Choice.
\end{proof}
\begin{definition}[Is Unique Matching Answer Choice]
  \label{def:physics:phyx-mini-0149:phyxminiproblems-problemphyxmini0149-isuniquematchinganswerchoice}
  \lean{PhyXMiniProblems.ProblemPhyXMini0149.IsUniqueMatchingAnswerChoice}
  \uses{def:physics:phyx-mini-0149:phyxminiproblems-problemphyxmini0149-focallengthestimatesetup, def:physics:phyx-mini-0149:phyxminiproblems-problemphyxmini0149-answerchoice, def:physics:phyx-mini-0149:phyxminiproblems-problemphyxmini0149-matchesanswerchoice}
  The selected choice is the unique displayed value equal to the estimate
\end{definition}
\begin{proof}
  This is the declared model definition of Is Unique Matching Answer Choice.
\end{proof}
\begin{lemma}[actual Focal Length is bracketed]
  \label{lem:physics:phyx-mini-0149:phyxminiproblems-problemphyxmini0149-actualfocallength-is-bracketed}
  \lean{PhyXMiniProblems.ProblemPhyXMini0149.actualFocalLength_is_bracketed}
  \uses{def:physics:phyx-mini-0149:phyxminiproblems-problemphyxmini0149-lengthincentimeters, def:physics:phyx-mini-0149:phyxminiproblems-problemphyxmini0149-focallengthestimatesetup, def:physics:phyx-mini-0149:phyxminiproblems-problemphyxmini0149-matchesproblemandprimaryfigure, def:physics:phyx-mini-0149:phyxminiproblems-problemphyxmini0149-satisfiesparaxialconverginglensmodel}
  The upright and inverted observations bracket the true focal length
\end{lemma}
\begin{proof}
  The conclusion follows from the governing relations and figure readouts named in the statement of actual Focal Length is bracketed.
\end{proof}
% --- Archon declaration topology end ---
% --- Archon physics formalization source end ---
